\chapter{Những Câu Nói Của Học Sinh}
\markboth{Những Câu Nói Của Học Sinh}{Things Students Often Say}

\section{``Em Không Hiểu.''}
\begin{truyen}{``Em Không Hiểu.''}{I Don't Understand}
\chuhoa{C}{âu nói này nghe rất nhỏ, nhưng nó là một cánh cửa lớn của việc học.}
\textit{(This sentence sounds small, but it opens a large door in learning.)}

Khi một học sinh dám nói mình chưa hiểu, em ấy đang thành thật với chính việc học của mình.
\textit{(When a student dares to say they do not understand, they are being honest with their own learning.)}

Điều đáng lo không phải là chưa hiểu, mà là không dám nói ra.
\textit{(The worrying thing is not confusion itself, but refusing to say it out loud.)}

Rất nhiều bài học bắt đầu tiến lên từ chính câu nói tưởng như yếu thế ấy.
\textit{(Many lessons begin to move forward from that seemingly weak sentence.)}
\end{truyen}

\begin{baihoc}
Biết nhận ra chỗ chưa hiểu là bước đầu của hiểu thật.
\textit{(Recognizing confusion is the first step toward real understanding.)}
\end{baihoc}

\begin{ghinhoanh}
\textbf{Short English to remember:}

Saying you do not understand is part of learning.

\end{ghinhoanh}

\begin{tuhocnhanh}
1. Vì sao câu ``em không hiểu'' lại quan trọng? \textit{Why is the sentence ``I don't understand'' important?}

2. Điều gì nguy hiểm hơn việc chưa hiểu? \textit{What is more dangerous than not understanding?}

3. Em có dám nói điều này khi cần không? \textit{Can you say this when needed?}
\end{tuhocnhanh}
\ngancach

\section{``Em Quên Làm Bài.''}
\begin{truyen}{``Em Quên Làm Bài.''}{I Forgot My Homework}
\chuhoa{N}{hiều học sinh nói câu này rất nhanh, như thể quên bài là chuyện nhỏ.}
\textit{(Many students say this quickly, as if forgetting homework were a small thing.)}

Nhưng phía sau chữ ``quên'' thường là một thói quen chưa nghiêm túc hoặc một cách sống thiếu sắp xếp.
\textit{(But behind the word ``forgot'' there is often weak discipline or a life lacking order.)}

Bài tập không chỉ để nộp. Nó là một phần của việc giữ nhịp học mỗi ngày.
\textit{(Homework is not only for submission. It keeps daily learning in rhythm.)}

Một lần quên có thể sửa được, nhưng quên mãi thì thành thói quen làm việc nửa chừng.
\textit{(One forgotten assignment can be fixed, but repeated forgetting becomes a habit of unfinished work.)}
\end{truyen}

\begin{baihoc}
Việc nhỏ làm không đến nơi thường báo trước những việc lớn cũng dễ dang dở.
\textit{(Small unfinished tasks often predict larger unfinished responsibilities.)}
\end{baihoc}

\begin{ghinhoanh}
\textbf{Short English to remember:}

Homework trains responsibility, not only memory.

\end{ghinhoanh}

\begin{tuhocnhanh}
1. Vì sao quên làm bài không chỉ là quên? \textit{Why is forgetting homework more than forgetting?}

2. Bài tập giúp rèn điều gì ngoài kiến thức? \textit{What does homework train besides knowledge?}

3. Em cần thay đổi gì để bớt quên? \textit{What do you need to change to forget less often?}
\end{tuhocnhanh}
\ngancach

\section{``Em Sợ Lên Bảng.''}
\begin{truyen}{``Em Sợ Lên Bảng.''}{I'm Afraid to Go to the Board}
\chuhoa{N}{ỗi sợ này rất thật với nhiều học sinh.}
\textit{(This fear is very real for many students.)}

Sợ sai, sợ bạn nhìn, sợ đứng trước lớp và thấy mình lúng túng.
\textit{(They fear mistakes, classmates' eyes, and the awkwardness of standing in front of everyone.)}

Nhưng lên bảng không chỉ là làm bài đúng. Nó còn là tập đứng trước sự run của chính mình.
\textit{(But going to the board is not only about being correct. It is also practice in standing inside one's own fear.)}

Có những học sinh tiến bộ không phải vì hết sợ, mà vì vẫn bước lên dù còn sợ.
\textit{(Some students improve not because fear disappears, but because they step forward anyway.)}
\end{truyen}

\begin{baihoc}
Dũng cảm trong lớp học thường bắt đầu bằng những bước rất run.
\textit{(Courage in the classroom often begins with trembling steps.)}
\end{baihoc}

\begin{ghinhoanh}
\textbf{Short English to remember:}

Courage often looks nervous at first.

\end{ghinhoanh}

\begin{tuhocnhanh}
1. Học sinh thường sợ điều gì khi lên bảng? \textit{What do students often fear at the board?}

2. Vì sao lên bảng không chỉ là chuyện đúng sai? \textit{Why is the board not only about right or wrong?}

3. Em có thể bước qua nỗi sợ này bằng cách nào? \textit{How can you step through this fear?}
\end{tuhocnhanh}
\ngancach

\section{``Em Làm Sai Rồi.''}
\begin{truyen}{``Em Làm Sai Rồi.''}{I Did It Wrong}
\chuhoa{C}{âu này có thể làm học sinh buồn, nhưng nó cũng rất đáng quý.}
\textit{(This sentence can make a student sad, yet it is also valuable.)}

Một người biết nhận ra mình sai thì vẫn còn cơ hội sửa.
\textit{(A person who can admit being wrong still has the chance to correct it.)}

Sai trong lúc học là chuyện bình thường hơn nhiều người tưởng.
\textit{(Being wrong while learning is far more normal than many think.)}

Điều quan trọng là từ chỗ sai đó, em nhìn ra mình cần học thêm điều gì.
\textit{(What matters is seeing what you need to learn from that mistake.)}
\end{truyen}

\begin{baihoc}
Nhận ra sai không làm em nhỏ đi. Nó làm em có cơ hội lớn lên.
\textit{(Recognizing error does not make you smaller. It gives you room to grow.)}
\end{baihoc}

\begin{ghinhoanh}
\textbf{Short English to remember:}

Admitting a mistake keeps growth possible.

\end{ghinhoanh}

\begin{tuhocnhanh}
1. Vì sao câu ``em làm sai rồi'' vẫn đáng quý? \textit{Why is ``I did it wrong'' still valuable?}

2. Sai trong lúc học có ý nghĩa gì? \textit{What does being wrong mean inside learning?}

3. Em thường phản ứng thế nào khi biết mình sai? \textit{How do you react when you realize you are wrong?}
\end{tuhocnhanh}
\ngancach

\section{``Thầy Giảng Lại Được Không?''}
\begin{truyen}{``Thầy Giảng Lại Được Không?''}{Could You Explain Again?}
\chuhoa{X}{in nghe lại một điều không phải dấu hiệu của kém cỏi.}
\textit{(Asking to hear something again is not a sign of weakness.)}

Nó cho thấy học sinh muốn theo kịp thay vì bỏ mặc chỗ hổng của mình.
\textit{(It shows a student wants to keep up instead of ignoring a gap.)}

Có nhiều em ngại hỏi lại vì sợ làm phiền hoặc sợ bị chê chậm.
\textit{(Many students hesitate to ask again because they fear being a burden or seeming slow.)}

Nhưng một câu hỏi đúng lúc có thể cứu cả một đoạn học tập phía sau.
\textit{(Yet one timely question can save an entire stretch of later learning.)}
\end{truyen}

\begin{baihoc}
Hỏi lại đúng lúc thường đỡ tốn hơn là im lặng sai rất lâu.
\textit{(Asking again at the right time often costs less than staying silently wrong for a long time.)}
\end{baihoc}

\begin{ghinhoanh}
\textbf{Short English to remember:}

Asking again can protect future learning.

\end{ghinhoanh}

\begin{tuhocnhanh}
1. Vì sao học sinh ngại xin giảng lại? \textit{Why do students hesitate to ask for another explanation?}

2. Một câu hỏi đúng lúc giúp gì cho việc học? \textit{How does a timely question help learning?}

3. Em có dám hỏi lại khi thật sự cần không? \textit{Can you ask again when you truly need it?}
\end{tuhocnhanh}
\ngancach

\section{``Bạn Làm Nhanh Quá.''}
\begin{truyen}{``Bạn Làm Nhanh Quá.''}{My Friend Is So Fast}
\chuhoa{S}{ự nhanh của người khác thường làm nhiều học sinh rối lòng.}
\textit{(Other people's speed often unsettles students.)}

Các em bắt đầu so mình với bạn và quên mất rằng mỗi người có nhịp hiểu khác nhau.
\textit{(They compare themselves and forget that everyone understands at a different pace.)}

Nhanh chưa chắc đã sâu, chậm cũng chưa chắc là dở.
\textit{(Fast does not always mean deep, and slow does not always mean poor.)}

Điều quan trọng là mình có đang tiến lên thật hay không.
\textit{(What matters is whether you are truly moving forward.)}
\end{truyen}

\begin{baihoc}
Đừng để tốc độ của người khác làm em quên nhịp học phù hợp của mình.
\textit{(Do not let another person's speed make you forget your own learning pace.)}
\end{baihoc}

\begin{ghinhoanh}
\textbf{Short English to remember:}

Learning pace is not the whole story.

\end{ghinhoanh}

\begin{tuhocnhanh}
1. Vì sao học sinh dễ bị áp lực vì bạn làm nhanh? \textit{Why do students feel pressured by faster classmates?}

2. Nhanh có luôn nghĩa là tốt hơn không? \textit{Does faster always mean better?}

3. Em cần giữ điều gì khi nhìn sang người khác? \textit{What should you hold onto when you look at others?}
\end{tuhocnhanh}
\ngancach

\section{``Em Thử Lại Nhé.''}
\begin{truyen}{``Em Thử Lại Nhé.''}{Let Me Try Again}
\chuhoa{Đ}{ây là một câu nói có sức sống rất đẹp trong lớp học.}
\textit{(This is a beautifully alive sentence in the classroom.)}

Nó cho thấy học sinh chưa muốn bỏ cuộc chỉ vì một lần chưa được.
\textit{(It shows a student does not want to quit after one unsuccessful attempt.)}

Thử lại không bảo đảm thành công ngay, nhưng nó giữ cho cánh cửa tiến bộ còn mở.
\textit{(Trying again does not guarantee immediate success, but it keeps the door of growth open.)}

Nhiều năng lực tốt được xây từ sự chịu thử lại như thế.
\textit{(Many strong abilities are built from that willingness to try again.)}
\end{truyen}

\begin{baihoc}
Người học bền thường không hơn người khác ở lần đầu, mà ở chỗ họ không bỏ sau lần đầu.
\textit{(Steady learners are often not better on the first try, but they do not stop after it.)}
\end{baihoc}

\begin{ghinhoanh}
\textbf{Short English to remember:}

Trying again keeps improvement alive.

\end{ghinhoanh}

\begin{tuhocnhanh}
1. Vì sao câu ``em thử lại nhé'' đáng quý? \textit{Why is ``let me try again'' valuable?}

2. Thử lại giữ lại điều gì cho việc học? \textit{What does trying again preserve in learning?}

3. Em có bỏ quá sớm ở đâu không? \textit{Where do you quit too early?}
\end{tuhocnhanh}
\ngancach

\section{``Em Hiểu Rồi.''}
\begin{truyen}{``Em Hiểu Rồi.''}{I Understand Now}
\chuhoa{Đ}{ó là một khoảnh khắc rất sáng trong lớp học.}
\textit{(It is one of the brightest moments in a classroom.)}

Có khi chỉ là một câu nói nhẹ, nhưng phía sau nó là cả một quãng loay hoay vừa được tháo ra.
\textit{(Sometimes it is just a soft sentence, but behind it lies a whole struggle finally untied.)}

Người dạy vui vì nhìn thấy học sinh chạm được vào điều trước đó còn mờ.
\textit{(The teacher is glad to see a student touch what was once unclear.)}

Người học cũng vui vì nhận ra mình có thể hiểu, chỉ là cần đúng cách và đủ thời gian.
\textit{(The learner is glad too, realizing understanding was possible with the right way and enough time.)}
\end{truyen}

\begin{baihoc}
Hiểu bài không chỉ là biết đáp án. Nó là cảm giác một chỗ mù trong đầu vừa sáng lên.
\textit{(Understanding is not only knowing the answer. It is the feeling of a dark spot in the mind lighting up.)}
\end{baihoc}

\begin{ghinhoanh}
\textbf{Short English to remember:}

Understanding often arrives after patience.

\end{ghinhoanh}

\begin{tuhocnhanh}
1. Vì sao khoảnh khắc ``em hiểu rồi'' đáng nhớ? \textit{Why is ``I understand now'' memorable?}

2. Nó cho thấy điều gì về việc học? \textit{What does it show about learning?}

3. Em nhớ lần nào mình thật sự hiểu ra một điều không? \textit{Do you remember a moment you truly understood something?}
\end{tuhocnhanh}
\ngancach

\section{``Hôm Nay Em Làm Được.''}
\begin{truyen}{``Hôm Nay Em Làm Được.''}{Today I Could Do It}
\chuhoa{N}{hiều tiến bộ trong lớp học đến từ những câu rất giản dị như vậy.}
\textit{(Much classroom progress appears in simple sentences like this.)}

Một bài làm được hôm nay có thể là điều em đã từng bó tay suốt mấy hôm trước.
\textit{(Something done today may be what felt impossible only days ago.)}

Niềm vui ấy nhỏ nhưng thật, và rất cần được giữ lại.
\textit{(That joy is small but real, and it deserves to be kept.)}

Người học tiến bộ thường không đi bằng những cú nhảy lớn, mà bằng nhiều ngày như thế.
\textit{(Growing learners often move not through giant leaps, but through many days like this.)}
\end{truyen}

\begin{baihoc}
Đừng coi thường một thành công nhỏ. Nhiều điều lớn được ghép từ những hôm ``mình làm được'' như thế.
\textit{(Do not underestimate a small success. Many big things are built from such days.)}
\end{baihoc}

\begin{ghinhoanh}
\textbf{Short English to remember:}

Small success deserves to be noticed.

\end{ghinhoanh}

\begin{tuhocnhanh}
1. Vì sao một lần làm được nhỏ vẫn đáng quý? \textit{Why is a small success still valuable?}

2. Tiến bộ thường đến theo kiểu nào? \textit{How does progress usually arrive?}

3. Hôm gần đây em đã làm được điều gì trước kia chưa làm được? \textit{What have you done recently that you could not do before?}
\end{tuhocnhanh}
\ngancach

\section{``Em Sẽ Cố Gắng.''}
\begin{truyen}{``Em Sẽ Cố Gắng.''}{I Will Keep Trying}
\chuhoa{Đ}{ây là một lời hứa rất quen của học sinh.}
\textit{(This is a very familiar promise from students.)}

Có lúc câu nói ấy chỉ thoáng qua, nhưng cũng có lúc nó là điểm bắt đầu của một thay đổi thật.
\textit{(Sometimes it passes quickly, but at other times it becomes the beginning of real change.)}

Cố gắng không làm mọi thứ dễ ngay, nhưng nó giữ cho con người không buông xuôi.
\textit{(Effort does not make everything easy at once, but it keeps people from giving up.)}

Trong lớp học, một lời hứa nhỏ với chính mình nhiều khi rất đáng được tin và nuôi dưỡng.
\textit{(In the classroom, a small promise to oneself is often worth trusting and nurturing.)}
\end{truyen}

\begin{baihoc}
Cố gắng không phải phép màu, nhưng nó là điều giữ đường học chưa bị khép lại.
\textit{(Effort is not magic, but it keeps the road of learning from closing.)}
\end{baihoc}

\begin{ghinhoanh}
\textbf{Short English to remember:}

Effort keeps the road open.

\end{ghinhoanh}

\begin{tuhocnhanh}
1. Vì sao câu ``em sẽ cố gắng'' vẫn quan trọng? \textit{Why is ``I will keep trying'' still important?}

2. Cố gắng giữ lại điều gì cho người học? \textit{What does effort preserve for a learner?}

3. Điều gì em thật sự muốn cố gắng hơn từ bây giờ? \textit{What do you truly want to work harder on from now?}
\end{tuhocnhanh}
\ngancach